\subsection{Environmental details}
\label{app:env}

\paragraph{Sim-to-sim tasks.}
We use four contact-rich manipulation tasks from MimicGen~\citep{mandlekar2023mimicgen}
--- Threading, Square, Coffee and ThreePieceAssembly --- and one locomotion task
from IsaacLab~\citep{mittal2023orbit}. All five are simulated, which is what lets
us run several hundred evaluation episodes per condition and, more importantly,
run the same controller with the \emph{true} target simulator in place of the
learned model (the oracle-dynamics arm of Sec.~\ref{sec:planner-adapt}). Neither
is possible on hardware.

Each task gives us a source domain and a target domain. The source domain is the
nominal simulator as shipped. The target domain is the same simulator with a
fixed perturbation applied to its physical parameters: we scale object mass and
inertia, the sliding friction of the manipulated object and of the fixture it
must be inserted into, the stiffness of the gripper actuator, and the geometric
clearance of the insertion. The task, the assets, the observation space and the
success predicate are untouched. This is a deliberate choice: because the two
domains differ only in physics, a policy that fails after transfer fails for a
dynamics reason and not because it is looking at a different scene or being
scored by a different rule.

Two consequences of that choice are worth stating explicitly. First, the success
predicate is a fixed geometric tolerance on the object pose and is \emph{not}
re-derived from the perturbed model, so shrinking an insertion clearance makes
the task physically harder without quietly redefining what counts as a success.
Second, the object's initial pose is drawn from the same distribution in both
domains, so the gap is not a perception or registration problem in disguise.

\paragraph{Target domain settings.}
Table~\ref{tab:app-targets} lists the perturbation for each task. Every entry is
a multiplicative scale on the corresponding field of the MuJoCo model, so
$1.0$ means unchanged. The scales were chosen per task to produce a large but
not insurmountable drop: the frozen source expert loses most of its success rate
in the target domain, while a policy trained directly in the target domain can
still solve the task, which is what makes the remaining gap attributable to
adaptation rather than to task difficulty.

\begin{table}[h]
\centering
\caption{Target-domain physical perturbations. Values are multiplicative scales
on the nominal simulator ($1.0$ = unchanged; \texttt{---} = not applicable to
that task). ``fixture friction'' and ``clearance'' refer to the part the object
is inserted into --- the ring for Threading, the peg for Square, the pod holder
for Coffee, and the base hole for ThreePieceAssembly. $\sigma$ is the standard
deviation of a Gaussian perturbation added to the robot's initial joint
configuration, applied identically in both domains.}
\label{tab:app-targets}
\small
\begin{tabular}{lcccc}
\toprule
 & Threading & Square & Coffee & ThreePiece \\
\midrule
object mass \& inertia      & $3.0$   & ---     & $8.0$   & ---     \\
object friction             & $0.5$   & ---     & $0.2$   & ---     \\
fixture friction            & $2.0$   & $2.0$   & $5.0$   & $2.0$   \\
insertion clearance         & $0.85$  & $1.2$\textsuperscript{$\dagger$} & --- & $0.96$ \\
gripper stiffness           & $0.12$  & $0.3$   & $0.3$   & $0.25$  \\
lid hinge damping           & ---     & ---     & $5.0$   & ---     \\
\midrule
episode length              & $400$   & $400$   & $400$   & $1000$  \\
initial-pose noise $\sigma$ & $0.1$   & $0.1$   & $0.1$   & $0$     \\
\bottomrule
\end{tabular}
\end{table}

\textsuperscript{$\dagger$}For Square the scale is applied to the peg rather
than to the hole, so a value \emph{above} one shrinks the clearance; the effect
is the same as the sub-one scales elsewhere in that row.

Square additionally restricts the square nut's initial yaw to $120$--$180^{\circ}$
in both domains. This is a narrower initial distribution than the default full
circle, and it is applied to the source and the target alike; it makes the task
tractable within our interaction budget without changing the gap between the two
domains. The other three tasks use the benchmark's default initial distribution.

% ---------------------------------------------------------------------------
% TODO(locomotion): the IsaacLab target is not fixed yet. When it is, add a
% short paragraph here rather than a fifth column -- the knobs are not the same
% quantities (no insertion clearance, no gripper) and forcing them into one
% table would invent correspondences that do not exist. State at minimum: which
% robot, which terrain/command distribution, what is perturbed in the target,
% the episode length, and whether the budget is counted in environment steps or
% in transitions (IsaacLab runs massively parallel environments, so "100k
% transitions" is not the same operating point as in the MimicGen tasks).
%
% TODO(consistency): ThreePieceAssembly currently runs with sigma = 0 while the
% other three use sigma = 0.1, and its episodes are 1000 steps rather than 400.
% Both are real differences in the setup, not typos -- either align them before
% the camera-ready or keep the row in the table as it is now. Do not silently
% report it as "the same protocol across four tasks".
%
% SOURCE OF TRUTH for every number above: the eval_domain_shift block of each
% task's SimDist stage-1 cfg.yaml, which is the env the adaptation runs actually
% instantiate (docs/infra.md 13). Threading:
% threading_simdist_dynens_diversev2_h5_seed1; Square:
% square_simdist_rotw1_dynens_h5_seed1; Coffee: coffee_s1g5x_dynens; TPA:
% tpa_simdist_diversev1_h5_seed1. Note Coffee has a second, older stage-1 with a
% DIFFERENT target (C50) whose directory name does not say so -- re-verify
% against stage1_ckpt, never against a directory name.
% ---------------------------------------------------------------------------
